% GENERATED FILE. Edit canonical lessons, then run npm run generate:books.
% canonical-source-hash: fnv1a64:bcacf8bca0146eb9
% canonical-lessons: ES-C380-simpatico, ES-C380-alegre, ES-C380-feo, ES-C380-ultimo, ES-C380-necesario

\chapter{Five Words for What Something Is Like}
\label{ch:simpatico-y-alegre}
\hlchaptermodality{380}

\begin{chapteropening}
\textbf{By the end of this chapter:} \emph{I can say simpático, alegre, feo, último and necesario, and count five people or things and say what each is like.}

Complete ES-C380-necesario: count five people or things and say what each is like, and say one thing about each.
\end{chapteropening}

\section[simpático]{simpático — a false friend worth knowing}

\label{lesson:ES-C380-simpatico}



\begin{quote}
Say \emph{lukewarm}, then \emph{thick}. (\emph{Tibio}; \emph{Grueso}.)
\end{quote}

\subsection*{What to know first}
\begin{itemize}
\raggedright
  \item amable — the nearest thing you already have.
  \item la gente — who you use it about.
\end{itemize}

\begin{sounds}
\begin{itemize}
\raggedright
  \item \emph{sim-PÁ-ti-co}, four beats with the weight on the second. Words stressed that far back always carry a written accent. $\to$ pronunciation reference
\end{itemize}
\end{sounds}

\begin{cousinweb}
\textbf{simpático} — \textquotedblleft{}likeable.\textquotedblright{}

\textbf{Simpático} means \textbf{likeable}. Pleasant company. Someone you are glad to sit next to.

It does \textbf{not} mean sympathetic, and this catches English speakers constantly. If you want sympathetic — feeling for somebody's trouble — Spanish reaches for \textbf{compasivo}.

The history explains the gap. Greek \textbf{sympátheia} meant a community of feeling; Latin borrowed it as \emph{sympathīa}; Spanish has it as \textbf{simpatía}. Then Spanish built its own adjective on top, and the Academy's dictionary gives \emph{simpático} no separate etymology at all — because there is nothing further back to give. The word was made in Spanish, and the sense \textquotedblleft{}likeable\textquotedblright{} is only about two and a half centuries old.

English \textbf{sympathetic} was built separately, on a different Greek suffix. So the two adjectives are cousins through the noun rather than the same word wearing two spellings, which is precisely why they were free to drift apart.

There is one place they still agree, and it is a strange one: the \emph{sistema nervioso simpático} is the sympathetic nervous system. Medicine kept the older meaning while ordinary speech walked off with it.
\end{cousinweb}

\begin{grammarlens}[title={What you've built}]
You can now describe somebody you like: \emph{muy simpático}.
\end{grammarlens}

\subsection*{Your turn}
Say these aloud:

\begin{itemize}
\raggedright
  \item \emph{simpático}
  \item \emph{muy simpático}
  \item \emph{una persona simpática}
\end{itemize}

\begin{tcolorbox}[breakable,colback=teal!4,colframe=teal!35!black,title={Before you move on}]
What does \emph{simpático} mean? (Likeable — not sympathetic.) And which Spanish word does mean sympathetic? (\emph{Compasivo}.)
\end{tcolorbox}
\section[alegre]{alegre — the word on a piece of sheet music}

\label{lesson:ES-C380-alegre}



\begin{quote}
Say \emph{thick}, then \emph{likeable}. (\emph{Grueso}; \emph{Simpático}.)
\end{quote}

\subsection*{What to know first}
\begin{itemize}
\raggedright
  \item contento — the quieter version.
  \item triste — the opposite.
\end{itemize}

\begin{sounds}
\begin{itemize}
\raggedright
  \item \emph{a-LE-gre}, three beats with the weight on the middle one. The \emph{g} before \emph{r} is hard, and the \emph{r} is a single tap. $\to$ pronunciation reference
\end{itemize}
\end{sounds}

\begin{cousinweb}
\textbf{alegre} — \textquotedblleft{}cheerful.\textquotedblright{}

Latin \textbf{alacer} meant lively, brisk, eager — quicker and more physical than \textquotedblleft{}cheerful.\textquotedblright{}

The Academy's dictionary does something worth noticing on the way to Spanish. Rather than going straight from \emph{alacer} to \emph{alegre}, it prints \textbf{two} reconstructed spoken forms in between, both marked with a star to say that nobody wrote them down. That is a dictionary telling you the road was bumpy and it is guessing at the potholes.

English has the word as \textbf{alacrity} — a brisk willingness, doing something without being asked twice. The old \textquotedblleft{}quick\textquotedblright{} sense survives there better than in Spanish.

And Italian handed it to music. \textbf{Allegro} on a score means fast, and every orchestra in the world reads a Latin word for liveliness at the top of the page.

Where \emph{alacer} came from before Latin is not settled. Several proposals exist; none has won.
\end{cousinweb}

\begin{grammarlens}[title={What you've built}]
You can now describe somebody in good spirits: \emph{muy alegre}.
\end{grammarlens}

\subsection*{Your turn}
Say these aloud:

\begin{itemize}
\raggedright
  \item \emph{alegre}
  \item \emph{muy alegre}
  \item \emph{una fiesta alegre}
\end{itemize}

\begin{tcolorbox}[breakable,colback=teal!4,colframe=teal!35!black,title={Before you move on}]
What did Latin \emph{alacer} mean? (Lively, brisk.) And which musical instruction is the same word? (\emph{Allegro}.)
\end{tcolorbox}
\section[feo]{feo — two Latin words that fell together}

\label{lesson:ES-C380-feo}



\begin{quote}
Say \emph{likeable}, then \emph{cheerful}. (\emph{Simpático}; \emph{Alegre}.)
\end{quote}

\subsection*{What to know first}
\begin{itemize}
\raggedright
  \item la cara — what it is usually said about.
  \item pálido — another word for how something looks.
\end{itemize}

\begin{sounds}
\begin{itemize}
\raggedright
  \item \emph{FE-o}, two beats with the weight on the first. The \emph{e} and the \emph{o} stay as two separate vowels — they do not run together. $\to$ pronunciation reference
\end{itemize}
\end{sounds}

\begin{cousinweb}
\textbf{feo} — \textquotedblleft{}ugly.\textquotedblright{}

Latin had \textbf{two} words spelled \emph{foedus}, and they had nothing to do with each other.

One was an \textbf{adjective}: foul, filthy, loathsome. That one is \emph{feo}.

The other was a \textbf{noun}: a treaty, a league, an alliance. That one is where English gets \textbf{federal}, \textbf{federation} and \textbf{confederate} — and, further back, it belongs with \emph{fidēs}, \textquotedblleft{}trust,\textquotedblright{} which Spanish keeps as \emph{fe} and \emph{fiel}. A federation is a thing built on trust.

The Academy's dictionary quietly keeps them apart. Under \emph{federal} it prints a gloss — \textquotedblleft{}pact, alliance.\textquotedblright{} Under \emph{feo} it prints none, because the adjective is a different word and needs no such explanation.

So \emph{feo} and \emph{federal} are one of the cleanest look-alikes in the language: a spelling collision inside Latin itself, carried forward into two modern languages that have no idea they are holding a coincidence.

One restraint worth keeping. It is safe to say the two Latin words are unrelated. It is less safe to give each of them a confident deeper ancestor — the treaty word's is well established, the ugly word's is not.
\end{cousinweb}

\begin{grammarlens}[title={What you've built}]
You can now describe something nobody wants to see: \emph{muy feo}.
\end{grammarlens}

\subsection*{Your turn}
Say these aloud:

\begin{itemize}
\raggedright
  \item \emph{feo}
  \item \emph{muy feo}
  \item \emph{un día feo}
\end{itemize}

\begin{tcolorbox}[breakable,colback=teal!4,colframe=teal!35!black,title={Before you move on}]
How many Latin words were spelled \emph{foedus}? (Two — an adjective meaning foul and a noun meaning treaty.) And which one gave English \emph{federal}? (The noun.)
\end{tcolorbox}
\section[último]{último — the far end of beyond}

\label{lesson:ES-C380-ultimo}



\begin{quote}
Say \emph{cheerful}, then \emph{ugly}. (\emph{Alegre}; \emph{Feo}.)
\end{quote}

\subsection*{What to know first}
\begin{itemize}
\raggedright
  \item primero — the other end.
  \item próximo — the one after this.
\end{itemize}

\begin{sounds}
\begin{itemize}
\raggedright
  \item \emph{ÚL-ti-mo}, three beats with the weight on the first. Words stressed that far back always carry a written accent. $\to$ pronunciation reference
\end{itemize}
\end{sounds}

\begin{cousinweb}
\textbf{último} — \textquotedblleft{}last.\textquotedblright{}

Latin \textbf{ultimus} meant farthest, last. It is a \textbf{superlative} — the furthest-out member of a family built on a stem meaning \textquotedblleft{}beyond.\textquotedblright{}

The family has three steps, and getting them in order matters:

\begin{itemize}
\raggedright
  \item the stem, meaning beyond;
  \item the comparative \textbf{ulter}, \textquotedblleft{}farther\textquotedblright{};
  \item the superlative \textbf{ultimus}, \textquotedblleft{}farthest.\textquotedblright{}
\end{itemize}

English borrowed from the end of that list. An \textbf{ultimate} answer is the last one there is. An \textbf{ultimatum} is the last offer before something unpleasant.

And \textbf{ultra} belongs to the family too — but not where people usually put it. \emph{Ultra} is not the base that \emph{ultimus} was built on. It is a frozen form of the comparative \emph{ulter}, which makes it \emph{ultimus}'s \textbf{sibling}, not its parent. Saying \textquotedblleft{}\emph{ultra} gave us \emph{ultimus}\textquotedblright{} reverses the relationship.

A small thing, and a good habit: when a family is laid out for you, check which member is the ancestor and which are the cousins standing beside it.
\end{cousinweb}

\begin{grammarlens}[title={What you've built}]
You can now name the final one: \emph{el último}.
\end{grammarlens}

\subsection*{Your turn}
Say these aloud:

\begin{itemize}
\raggedright
  \item \emph{último}
  \item \emph{el último}
  \item \emph{la última vez}
\end{itemize}

\begin{tcolorbox}[breakable,colback=teal!4,colframe=teal!35!black,title={Before you move on}]
What kind of form is \emph{ultimus}? (A superlative — the farthest.) And is \emph{ultra} its parent? (No — its sibling, a frozen form of the comparative.)
\end{tcolorbox}
\section[necesario]{necesario — a word two languages read out of the same book}

\label{lesson:ES-C380-necesario}



\begin{quote}
Say \emph{ugly}, then \emph{last}. (\emph{Feo}; \emph{Último}.)
\end{quote}

\subsection*{What to know first}
\begin{itemize}
\raggedright
  \item bastante — the word for having enough.
  \item poco — the word for not having enough.
\end{itemize}

\begin{sounds}
\begin{itemize}
\raggedright
  \item \emph{ne-ce-SA-rio}, four beats with the weight on the third. The \emph{c} before \emph{e} is soft, and the \emph{io} at the end is one glide. $\to$ pronunciation reference
\end{itemize}
\end{sounds}

\begin{cousinweb}
\textbf{necesario} — \textquotedblleft{}necessary.\textquotedblright{}

Latin \textbf{necessarius} came from \textbf{necesse}, \textquotedblleft{}unavoidable.\textquotedblright{} Spanish has \emph{necesario}; English has \textbf{necessary}.

Neither language inherited it. Both \textbf{read it out of Latin} — Spanish and English each lifted the word from the written language rather than receiving it through centuries of speech. That is why they still match almost letter for letter:

\begin{itemize}
\raggedright
  \item \emph{necesario} and \emph{necessary};
  \item \emph{necesidad} and \emph{necessity}.
\end{itemize}

Compare that with a word that was actually spoken the whole way, where the two languages end up unrecognisable to each other. A pair that still line up this neatly are usually a pair that were \textbf{borrowed}, not inherited, and that is a reliable thing to notice.

On where \emph{necesse} itself came from: nobody has closed the case. There is a defensible proposal, and there is a tempting gloss attached to it — that \emph{necesse} means \textquotedblleft{}what cannot be got round.\textquotedblright{} Set that aside. The gloss is the proposal restated and nothing more: you get \textquotedblleft{}cannot be got round\textquotedblright{} only by first assuming the analysis it is supposed to support. A translation of a guess is not evidence for the guess.
\end{cousinweb}

\begin{grammarlens}[title={What you've built}]
You can now say that something must happen: \emph{es necesario}.
\end{grammarlens}

\subsection*{Your turn}
Say these aloud:

\begin{itemize}
\raggedright
  \item \emph{necesario}
  \item \emph{es necesario}
  \item \emph{no es necesario}
\end{itemize}

\begin{tcolorbox}[breakable,colback=teal!4,colframe=teal!35!black,title={Before you move on}]
Did Spanish inherit \emph{necesario} or borrow it? (Borrow it, out of written Latin — which is why English matches it so closely.) And is it known what \emph{necesse} is built from? (No; the case is open.)
\end{tcolorbox}
